\chapter{Twierdzenie Taylora}
\label{ma1:ch:taylor}

Pochodna daje przybliżenie liniowe funkcji w otoczeniu punktu. Pochodne wyższych rzędów pozwalają zastąpić funkcję wielomianem z większą dokładnością. Rozróżnimy przybliżenie lokalne, opisane resztą Peana, oraz oszacowania błędu na całym odcinku, wynikające z twierdzenia o wartości średniej.

\section{Wielomian Taylora}
\begin{definicja}
\label{ma1:def:taylor-polynomial}
Niech $n\ge0$ będzie liczbą całkowitą, a $f:I\to\R$ funkcją na otwartym przedziale $I$. Jeśli pochodne do rzędu $n$ istnieją w punkcie $a\in I$ w sensie definicji \ref{ma1:def:higher-derivatives}, to wielomianem Taylora stopnia co najwyżej $n$ funkcji $f$ w punkcie $a$ nazywamy
\[
 T_{n,a}f(x)=\sum_{k=0}^n\frac{f^{(k)}(a)}{k!}(x-a)^k.
\]
Resztą jest $R_{n,a}f(x)=f(x)-T_{n,a}f(x)$ dla $x\in I$. Gdy $a=0$, mówi się także o wielomianie Maclaurina. Przyjmujemy $(x-a)^0=1$ oraz $0!=1$.
\end{definicja}
\begin{stwierdzenie}
\label{ma1:prop:taylor-jet}
Wielomian $T_{n,a}f$ jest jedynym wielomianem $P$ stopnia co najwyżej $n$ spełniającym $P^{(j)}(a)=f^{(j)}(a)$ dla $0\le j\le n$.
\end{stwierdzenie}
\begin{proof}
Każdy taki wielomian ma jednoznaczny zapis $P(x)=\sum_{k=0}^n c_k(x-a)^k$, uzyskany przez zapis wielomianu $P(a+t)$ w potęgach $t$. Dla $j>k$ pochodna $j$-tego rzędu potęgi jest zerowa, dla $j=k$ wynosi $k!$, a dla $j<k$ zawiera czynnik $(x-a)^{k-j}$, który znika w $a$ (stwierdzenie \ref{ma1:prop:higher-elementary} i reguła łańcuchowa). Stąd $P^{(j)}(a)=j!c_j$. Warunki wyznaczają jednoznacznie $c_j=f^{(j)}(a)/j!$, a te współczynniki rzeczywiście spełniają wszystkie równości.
\end{proof}

\section{Postać Peana}
Zapis $r(x)=o((x-a)^n)$ przy $x\to a$ oznacza w tym rozdziale
$\lim_{x\to a,\,x\ne a}r(x)/(x-a)^n=0$; dla $n=0$ mianownik wynosi $1$.

\begin{twierdzenie}[Taylora z resztą Peana]
\label{ma1:thm:taylor-peano}
Niech $a\in I$, gdzie $I$ jest otwartym przedziałem. Dla $n=0$ załóżmy ciągłość $f$ w $a$. Dla $n\ge1$ załóżmy, że $f$ jest $(n-1)$ razy różniczkowalna w otoczeniu $a$ i że $f^{(n-1)}$ jest różniczkowalna w $a$. Wówczas
\[
 f(x)=T_{n,a}f(x)+o((x-a)^n)\qquad(x\to a).
\]
Nie zakładamy ciągłości $f^{(n)}$ ani nawet jej istnienia w innych punktach otoczenia.
\end{twierdzenie}
\begin{proof}
Dla $n=0$ teza jest ciągłością w $a$. Dla $n=1$ wynika wprost z definicji pochodnej:
\[
 \frac{f(x)-f(a)-f'(a)(x-a)}{x-a}
 =\frac{f(x)-f(a)}{x-a}-f'(a)\longrightarrow0.
\]
Niech $n\ge2$ i załóżmy tezę dla rzędu $n-1$. Funkcja $f'$ spełnia założenia tej tezy: jej pochodne do rzędu $n-2$ istnieją w otoczeniu $a$, a $(f')^{(n-2)}=f^{(n-1)}$ jest różniczkowalna w $a$. Dla $R=f-T_{n,a}f$ mamy $R(a)=0$ i
\[
 R'(t)=f'(t)-\sum_{j=0}^{n-1}\frac{f^{(j+1)}(a)}{j!}(t-a)^j
       =o((t-a)^{n-1})
\]
przy $t\to a$, na mocy założenia indukcyjnego. Dla $x\ne a$ dostatecznie bliskiego $a$ funkcje $R$ i $g(t)=(t-a)^n$ są ciągłe na odcinku o końcach $a,x$ i różniczkowalne w jego wnętrzu; ciągłość $R$ wynika z różniczkowalności $f$ i wielomianu. Ponadto $g'(t)=n(t-a)^{n-1}\ne0$ we wnętrzu. Wniosek \ref{ma1:cor:cauchy-ratio}, z uporządkowanymi końcami odcinka (zmiana ich kolejności zmienia znak obu przyrostów), daje punkt $c_x$ ściśle między $a$ i $x$ taki, że
\[
 \frac{R(x)}{(x-a)^n}=\frac{R'(c_x)}{n(c_x-a)^{n-1}}.
\]
Mamy $0<|c_x-a|<|x-a|$, więc $c_x\to a$ przy $x\to a$ i prawa strona dąży do zera. To kończy indukcję.
\end{proof}
\begin{wniosek}
\label{ma1:cor:taylor-peano-unique}
Przy założeniach twierdzenia \ref{ma1:thm:taylor-peano} wielomian Taylora jest jedynym wielomianem $P$ stopnia co najwyżej $n$, dla którego $f(x)-P(x)=o((x-a)^n)$.
\end{wniosek}
\begin{proof}
Jeśli $P$ ma tę własność, to $Q=P-T_{n,a}f=o((x-a)^n)$ przez odjęcie dwóch reszt. Gdyby $Q\ne0$, istniałby najmniejszy indeks $j\le n$ o niezerowym współczynniku $d_j$ w rozwinięciu $Q(x)=\sum_{k=j}^n d_k(x-a)^k$. Wtedy $Q(x)/(x-a)^j\to d_j\ne0$. Z drugiej strony iloraz ten jest równy $(Q(x)/(x-a)^n)(x-a)^{n-j}$, co dąży do zera: pierwszy czynnik dąży do zera, drugi jest ograniczony w otoczeniu $a$. Sprzeczność daje $Q=0$.
\end{proof}

\section{Postacie reszty wynikające z twierdzenia Cauchy'ego}
\begin{twierdzenie}
\label{ma1:thm:taylor-remainders}
Niech $a,x\in\R$, $a\ne x$, $n\ge0$ będzie liczbą całkowitą, a $K=[\min(a,x),\max(a,x)]$. Niech $f\in C^n(K)$ w sensie rozdziału \ref{ma1:ch:higher-derivatives} i niech $f^{(n+1)}$ istnieje we wnętrzu $K$. Dla każdego całkowitego $p\ge1$ istnieje $\xi$ ściśle między $a$ i $x$, dla którego
\[
 R_{n,a}f(x)=
 \frac{f^{(n+1)}(\xi)}{p\,n!}
 (x-\xi)^{n+1-p}(x-a)^p.
\]
Pochodne na końcach rozumiemy jako ciągłe przedłużenia pochodnych z wnętrza. W szczególności dla $p=n+1$ otrzymujemy resztę Lagrange'a, a dla $p=1$ resztę Cauchy'ego:
\[
 \begin{aligned}
 R_{n,a}f(x)&=\frac{f^{(n+1)}(\xi_L)}{(n+1)!}(x-a)^{n+1},\\
 R_{n,a}f(x)&=\frac{f^{(n+1)}(\xi_C)}{n!}(x-\xi_C)^n(x-a).
 \end{aligned}
\]
Punkty $\xi_L,\xi_C$ nie muszą być równe.
\end{twierdzenie}
\begin{proof}
Ustalmy $x$ i rozważmy na $K$ funkcje
\[
 F(t)=f(x)-\sum_{k=0}^n\frac{f^{(k)}(t)}{k!}(x-t)^k,
 \qquad G(t)=(x-t)^p.
\]
Ciągłość pochodnych do rzędu $n$ zapewnia ciągłość $F$ na $K$; istnienie pochodnych do rzędu $n+1$ we wnętrzu zapewnia jej różniczkowalność tam. Funkcja $G$ jest wielomianem. Mamy $F(a)=R_{n,a}f(x)$, $F(x)=0$, $G(a)=(x-a)^p\ne0$ i $G(x)=0$. Dla $t$ we wnętrzu odcinka, różniczkując każdy składnik, dostajemy
\[
 F'(t)=-\sum_{k=0}^n\frac{f^{(k+1)}(t)}{k!}(x-t)^k
       +\sum_{k=1}^n\frac{f^{(k)}(t)}{(k-1)!}(x-t)^{k-1}.
\]
Dla $n=0$ druga suma jest pusta. Po podstawieniu $j=k-1$ w drugiej sumie jej składniki kasują dokładnie składniki pierwszej o indeksach $0,\ldots,n-1$. Pozostaje
\[
 F'(t)=-\frac{f^{(n+1)}(t)}{n!}(x-t)^n,
 \qquad G'(t)=-p(x-t)^{p-1}\ne0.
\]
Wniosek \ref{ma1:cor:cauchy-ratio} zastosowany do $F,G$ na $K$ daje punkt $\xi$ w jego wnętrzu z
\[
 \frac{R_{n,a}f(x)}{(x-a)^p}
 =\frac{F'(\xi)}{G'(\xi)}
 =\frac{f^{(n+1)}(\xi)}{p\,n!}(x-\xi)^{n+1-p}.
\]
Równość jest niezależna od kolejności końców, ponieważ oba przyrosty zmieniają wtedy znak. Mnożymy przez $(x-a)^p$ i podstawiamy kolejno $p=n+1$ oraz $p=1$. W przypadku ujemnego wykładnika w ogólnym wzorze nie ma dzielenia przez zero, gdyż $\xi\ne x$.
\end{proof}

\begin{wniosek}[Oszacowanie reszty]
\label{ma1:cor:taylor-error-bound}
Jeżeli przy założeniach twierdzenia \ref{ma1:thm:taylor-remainders} zachodzi $|f^{(n+1)}(t)|\le M$ we wnętrzu $K$, to
\[
 |R_{n,a}f(x)|\le \frac{M}{(n+1)!}|x-a|^{n+1}.
\]
Jeżeli $x=a$, reszta jest równa zeru.
\end{wniosek}
\begin{proof}
Bierzemy moduły obu stron wzoru Lagrange'a i korzystamy z ograniczenia pochodnej w punkcie $\xi_L$. Dla $x=a$ jedynym niezerowym potencjalnie składnikiem wielomianu jest $f(a)$, więc $R_{n,a}f(a)=0$.
\end{proof}
W zapisie $\xi=a+\theta(x-a)$ parametr $\theta=(\xi-a)/(x-a)$ należy do $(0,1)$ niezależnie od znaku $x-a$. Postać całkową reszty wyprowadzimy dopiero po wprowadzeniu całki.

\section{Przykłady i kontrola dokładności}
\begin{przyklad}
\label{ma1:ex:taylor-exp}
Z wzorów na pochodne funkcji wykładniczej (stwierdzenie \ref{ma1:prop:higher-elementary}) dla $f(x)=e^x$ wynika $f^{(k)}(0)=1$, więc $T_{n,0}f(x)=\sum_{k=0}^n x^k/k!$. Dla $|x|\le r$, $r>0$, wszystkie punkty odcinka od $0$ do $x$ należą do $[-r,r]$, gdzie $e^t\le e^r$. Wniosek \ref{ma1:cor:taylor-error-bound} daje
\[
 \left|e^x-\sum_{k=0}^n\frac{x^k}{k!}\right|
 \le e^r\frac{r^{n+1}}{(n+1)!}.
\]
Prawa strona dąży do zera dla ustalonego $r$, ponieważ iloraz kolejnych wyrazów $r^{n+1}/(n+1)!$ wynosi $r/(n+2)$, a od pewnego indeksu jest nie większy niż $1/2$. Dalsze wyrazy są więc ograniczone przez ciąg geometryczny dążący do zera.
\end{przyklad}
\begin{przyklad}
\label{ma1:ex:taylor-trig}
Dla $\sin$ i $\cos$ pochodne w zerze wynikają z cyklu czterech pochodnych ze stwierdzenia \ref{ma1:prop:higher-elementary}. Dla $m\ge0$:
\[
 T_{2m+2,0}\sin(x)=\sum_{j=0}^{m}\frac{(-1)^j x^{2j+1}}{(2j+1)!},
 \quad
 T_{2m+1,0}\cos(x)=\sum_{j=0}^{m}\frac{(-1)^j x^{2j}}{(2j)!}.
\]
Każda pochodna obu funkcji ma moduł nie większy niż $1$. Z wniosku \ref{ma1:cor:taylor-error-bound} błędy powyższych przybliżeń są ograniczone odpowiednio przez $|x|^{2m+3}/(2m+3)!$ i $|x|^{2m+2}/(2m+2)!$. Wykorzystaliśmy zerowe współczynniki kolejnych rzędów: wielomiany nie zmieniają się po dopisaniu tych wyrazów, lecz oszacowanie reszty jest wyższego rzędu.
\end{przyklad}
\begin{przyklad}
\label{ma1:ex:taylor-log}
Dla $f(x)=\log(1+x)$ na $(-1,\infty)$ wzory na pochodne logarytmu i reguła łańcuchowa dają $f^{(k)}(0)=(-1)^{k-1}(k-1)!$ dla $k\ge1$, a $f(0)=0$. Dla $n\ge1$,
\[
 T_{n,0}f(x)=\sum_{k=1}^n\frac{(-1)^{k-1}x^k}{k}.
\]
Jeśli $x\ge0$, to we wnętrzu odcinka od $0$ do $x$ mamy $1+t\ge1$ i $|f^{(n+1)}(t)|=n!/(1+t)^{n+1}\le n!$. Stąd błąd nie przekracza $x^{n+1}/(n+1)$. Dla $-1<x<0$ w tym odcinku mamy $1+t\ge1+x>0$, więc analogiczne oszacowanie wynosi
\[
 |R_{n,0}f(x)|\le\frac{1}{n+1}
                 \left(\frac{|x|}{1+x}\right)^{n+1}.
\]
Ostatnie oszacowanie nie gwarantuje zaniku błędu dla wszystkich $-1<x<0$; nie należy wyciągać z niego takiego wniosku. Dokładniejszą tożsamość można otrzymać z reszty Cauchy'ego: dla $\xi\in(x,0)$ mamy $0<(\xi-x)/(1+\xi)\le |x|$ (po pomnożeniu przez dodatnie $1+\xi$ nierówność sprowadza się do $\xi(1+x)\le0$). Zatem
\[
 |R_{n,0}f(x)|
 =\frac{|x|}{1+\xi}\left(\frac{\xi-x}{1+\xi}\right)^n
 \le \frac{|x|^{n+1}}{1+x}.
\]
Ten błąd dąży do zera dla każdego ustalonego $-1<x<0$.
\end{przyklad}

\begin{center}
\begin{tikzpicture}[x=4.6cm,y=1.2cm,>=Stealth]
 \draw[->] (-.7,0)--(1.32,0) node[right] {$x$};
 \draw[->] (0,-.12)--(0,5.8) node[above] {$y$};
 \draw[blue!65!black,thick,domain=-.55:1.12,samples=100]
 plot (\x,{1+\x+\x*\x+\x*\x*\x});
 \draw[orange!85!black,dashed,thick,domain=-.55:1.12,samples=100]
 plot (\x,{1+\x+\x*\x});
 \draw[gray,dotted] (1,0)--(1,3);
 \draw[<->,green!40!black,thick] (1,3)--(1,4)
 node[midway,right,font=\small] {$R_{2,0}f(1)$};
 \fill (0,1) circle (1.5pt);
 \fill[blue!65!black] (1,4) circle (1.5pt);
 \fill[orange!85!black] (1,3) circle (1.5pt);
 \node[below left] at (0,0) {$a=0$};
 \node[below] at (1,0) {$1$};
 \node[blue!65!black,above left,font=\small] at (.88,3.335872) {$f$};
 \node[orange!85!black,below right,font=\small] at (.7,2.19) {$T_{2,0}f$};
\end{tikzpicture}

\small Rysunek: dla $f(x)=1+x+x^2+x^3$ wielomian $T_{2,0}f(x)=1+x+x^2$ ma w zerze te same wartości pochodnych do rzędu drugiego; w punkcie $1$ reszta wynosi $1$.
\end{center}

\section{Gładkość a odtwarzanie funkcji przez wielomiany}
\begin{przyklad}
\label{ma1:ex:taylor-flat}
Niech $H(0)=0$ oraz $H(x)=e^{-1/x^2}$ dla $x\ne0$. Funkcja jest klasy $C^\infty(\R)$, a wszystkie jej pochodne w zerze są równe zeru, mimo że $H(x)>0$ dla $x\ne0$.

Najpierw dla każdej całkowitej $N\ge0$ pokażemy, że $|x|^{-N}e^{-1/x^2}\to0$ przy $x\to0$, $x\ne0$. Wybierzmy całkowite $k>N/2$. Dla $u>0$ wzór Lagrange'a dla $e^u$ wokół zera daje $e^u\ge u^k/k!$, ponieważ wszystkie składniki wielomianu oraz reszta są dodatnie. Dla $u=1/x^2$ otrzymujemy $0\le |x|^{-N}e^{-1/x^2}\le k!|x|^{2k-N}\to0$.

Poza zerem pochodne mają postać $H^{(j)}(x)=P_j(1/x)e^{-1/x^2}$, gdzie $P_0=1$ i
\[
 P_{j+1}(z)=-z^2P_j'(z)+2z^3P_j(z).
\]
Równość wynika z reguł łańcuchowej i iloczynu, gdyż $(1/x)'=-1/x^2$ oraz $(-1/x^2)'=2/x^3$. Każdy $P_j$ jest wielomianem. Dla $|x|\le1$, $x\ne0$, jego wartość w $1/x$ jest ograniczona przez $C|x|^{-N}$, gdzie $N$ jest stopniem wielomianu i $C$ sumą modułów współczynników (dla wielomianu zerowego przyjmujemy $C=N=0$).

Indukcyjnie, jeśli $H^{(j)}(0)=0$, to iloraz różnicowy w zerze wynosi $P_j(1/h)e^{-1/h^2}/h$ i jego moduł nie przekracza $C|h|^{-N-1}e^{-1/h^2}\to0$. Stąd $H^{(j+1)}(0)=0$. Przypadek początkowy to definicja $H(0)$. Ten sam argument z wykładnikiem $N$, bez dzielenia przez $h$, dowodzi ciągłości każdej pochodnej w zerze; poza zerem ciągłość wynika z podanej postaci. Zatem $H\in C^\infty(\R)$. Wszystkie jej wielomiany Taylora w zerze są zerowe i nie odtwarzają funkcji w żadnym otoczeniu zera. Gładkość sama w sobie nie wystarcza do odtwarzania funkcji przez jej rozwinięcie Taylora.
\end{przyklad}
